\documentclass[10pt,a4paper]{amsart}
\usepackage[margin=19mm]{geometry}
\usepackage{amsmath,amssymb,mathtools}
\usepackage[scr=rsfs,cal=euler]{mathalfa}
\usepackage{xfrac}
\usepackage[unicode,hidelinks,bookmarks=false]{hyperref}
\newcommand{\ie}{i.e.\ }
\newcommand{\cf}{cf.\ }
\newcommand{\resp}{respectively\ }
\newcommand{\Subsection}{\subsection}
\providecommand{\nobreakdash}{\nobreak\hbox{-}}
\setlength{\emergencystretch}{2em}
\multlinegap0pt
\usepackage{bm}
\usepackage{bbm}
\usepackage{dsfont}
\usepackage{xspace}
\usepackage{cancel}
\usepackage{mathtools}
\usepackage{amsthm}
\makeatletter
\@ifundefined{theorem}{\newtheorem{theorem}{Theorem}}{}
\@ifundefined{definition}{\newtheorem{definition}[theorem]{Definition}}{}
\@ifundefined{lemma}{\newtheorem{lemma}[theorem]{Lemma}}{}
\makeatother
\title[On the graded singularity category of Abelian quotient singu]{On the graded singularity category of Abelian quotient singularities, I. Smooth categorical compactification}
\author{Xiaojun Chen, Jieheng Zeng}
\date{Source: arXiv:2507.19815v1 (CC BY 4.0)}
\begin{document}
\maketitle

\noindent\emph{Source and licence.} On the graded singularity category of Abelian quotient singularities, I. Smooth categorical compactification. Version arXiv:2507.19815v1 (CC BY 4.0), distributed under the Creative Commons Attribution 4.0 International licence, as recorded in the arXiv licence metadata for this exact version and shown on the abstract page https://arxiv.org/abs/2507.19815v1. Full rights notice: licenses/arxiv-2507.19815-CC-BY-4.0.txt.\par\smallskip

\noindent\textit{Editorial scope.} Counted unit: the Lemma whose source label is \texttt{NNned1} in the article (source0.tex lines 2602--2604, source label \texttt{NNned1}), delivered together with the article's own complete original proof of it (source0.tex lines 2606--2881, one \texttt{proof} environment of 276 lines). No mathematics is written, repaired or re-derived: statement and proof are the article's own text line for line. Required context retained uncounted: Inject (lines 2120--2127); Compos (lines 2244--2251); bdfbv11267 (lines 2563--2571); jxkdbn (lines 2594--2597). Every definition, display and label of those lines is retained unchanged. Omitted from this package and not redistributed: 1--2119, 2128--2243, 2252--2562, 2572--2593, 2598--2601, 2605--2605, 2882--4375. Nothing inside a retained excerpt is omitted or altered: every retained body is delivered exactly as the article's lines are written, so no omission receipt of any kind accompanies this package. Citations: the retained excerpts carry no citation command at all, so no bibliography entry of the article is transcribed and no reference is redistributed. Reference adaptation: none was needed and none was made; every reference inside the delivered unit resolves inside it, so no unresolved reference or citation is produced. Printed slips kept literal: no printed word, symbol, hypothesis or number of the counted statement or of its proof was altered. Layout and class: the article's own class is \texttt{amsart} with packages including geometry, graphicx, epstopdf, subfigure, latexsym, bm, amsmath, amsthm, amssymb, amsfonts, eucal, xy, mathrsfs, color; the wrapper is the lane's pinned \texttt{amsart} editorial wrapper at 10pt on A4, which restates the theorem-like environments the retained text needs and adds the article's own macro definitions that the retained text needs (none), with extra wrapper packages (bm, bbm, dsfont, xspace, cancel, mathtools). The delivered numbering is the unit's own. Assets: no retained span contains a figure environment, a graphics-inclusion command, a PDF-page inclusion, a table or a TikZ picture; no image file is copied into this package, the package declares no asset dependency and no image file is redistributed with it. Licence: version arXiv:2507.19815v1 of this article is distributed under the Creative Commons Attribution 4.0 International licence, as recorded in the arXiv licence metadata for this exact version and shown on the abstract page https://arxiv.org/abs/2507.19815v1. A full rights notice is delivered at licenses/arxiv-2507.19815-CC-BY-4.0.txt. Characters: every retained excerpt is pure ASCII, so the delivered engine, XeLaTeX with its default Latin Modern text font, reports no missing character.
\begin{lemma}\label{Inject}
Fixing character $\chi$
and
its corresponding indecomposable idempotent $e_{\chi}$, there is a graded algebra injection 
$$
(G / H^j) \sharp k[W^{j}] \hookrightarrow G \sharp k[V] \cong G \sharp R.
$$		
\end{lemma}

\begin{lemma}\label{Compos}
With the above settings, there is a graded algebra surjection
$$
	\bar{\varsigma}^{j}_{\chi}: G \sharp R  \twoheadrightarrow   (G / H^j) \sharp k[W^{j}]
$$
such that $\bar\varsigma^{j}_{\chi} \circ \bar\kappa^{j}_{\chi} 
= \mathrm{Id}$ in $(G / H^j) \sharp k[W^{j}]$. 
\end{lemma}

\begin{definition}\label{bdfbv11267}
Let
$$
\sqrt{\Lambda e\Lambda} := \mathrm{Span}_{k}\Big\{e' \otimes f 
\in \Lambda \big\vert \exists m \in \mathbb{N}
\mbox{ such that }
e' \otimes f^m \in \Lambda e\Lambda \Big\}. 
$$
\end{definition}

\begin{align}\label{jxkdbn}
\bar{\varpi}_0 : =  \bigoplus_{(\chi, G^i) \in \tilde{G}_0} \bar{\eta}^{i}_{\chi}: 
\Lambda / \sqrt{\Lambda e \Lambda} \rightarrow \bigoplus_{(\chi, G^i) \in \tilde{G}_0} \Lambda_i.
\end{align}

\begin{lemma}\label{NNned1}
$\bar{\varpi}_0$ is injective. 
\end{lemma}

\begin{proof}
This proof is a bit involved and
consists of several steps.

\vspace{2mm}

\noindent{\bf Step 1. We
reduce the
proof of the lemma
to the proof of the following (\ref{hxowc453}).}

In fact, by the definition 
(\ref{jxkdbn}), notice that to 
prove the lemma,
it suffices to show that 
$$
\mathrm{Ker}(\bar{\varpi}_0) \cong \bigcap\limits_{(\chi, G^i) \in \tilde{G}_0}
 \mathrm{Ker}(\bar{\eta}^{i}_{\chi})	
$$	
is trivial on $\Lambda / \sqrt{\Lambda e \Lambda}$. 


Since $\eta^{i}_{\chi}$ is also a graded $\Lambda$-module morphism,
$\mathrm{Ker}(\eta^{i}_{\chi})$ is a graded $\Lambda$-module. Thus, it suffices to show that for any 
indecomposable idempotent $u$ of $\Lambda$, 
\begin{align}\label{qsxbd83}	
\Big(\bigcap\limits_{(\chi, G^i) \in \tilde{G}_0} \mathrm{Ker}(\eta^{i}_{\chi}) \Big)\otimes_{\Lambda} \Lambda u 
& \cong  \bigcap\limits_{(\chi, G^i) \in \tilde{G}_0} ( \mathrm{Ker}(\eta^{i}_{\chi}) 
\otimes_{\Lambda} \Lambda u ) \nonumber \\
 & \cong \bigcap\limits_{(\chi, G^i) \in \tilde{G}_0} \mathrm{Ker} (\eta^{i}_{\chi} 
 \otimes_{\Lambda} \Lambda u ) 
\end{align}
is trivial as graded $R$-modules. 

Next, $\sqrt{\Lambda e \Lambda} \otimes_{\Lambda} \Lambda u \cong (\sqrt{\Lambda e \Lambda} )u
\subseteq \Lambda u \cong R$ is a reduced ideal, denoted by $I_u$, of $R$.  
%
Let $\chi_u$ be the character of $G$ corresponding to $u$. 
Then we have graded algebra homomorphism $$\eta^{i}_{\chi} \otimes_{\Lambda} \Lambda u : 
R/I_u \cong \Lambda u \big/ (\sqrt{\Lambda e \Lambda} )u \twoheadrightarrow  \Lambda_i \otimes_{\Lambda} \Lambda u \cong k[V^i].$$
It is nontrivial if and only 
if $(\chi_u, G^i) \sim (\chi, G^i)$ in $\tilde{G}_0$. 
Now assume
that $(\chi_u, G^i) \sim (\chi, G^i)$. 
Let $\iota_{\#}^i : R \twoheadrightarrow k[V^i]$ be the canonical projection and $I^i$ be its kernel. 
Here, we note that
$\iota_{\#}^i$ induces $\eta^{i}_{\chi} 
\otimes_{\Lambda} \Lambda u : R/ I_u \twoheadrightarrow R/I^i$. 

Now, by the construction of $\tilde{G}_0$, to prove (\ref{qsxbd83}) is trivial, it is enough to show that 
\begin{align}\label{2736sdd}
\bigcap\limits_{(\chi_0, G^i) \not\sim (\chi_u, G^i)} I^i \subseteq I_u.
\end{align}
To this end, assume 
$f \in \bigcap\limits_{(\chi_0, G^i) 
\not\sim (\chi_u, G^i)} I^i $ is 
a monomial.
We need to show $f \in I_u$. 
In fact,
by the definition of 
$\sqrt{\Lambda e \Lambda}$ 
(see Definition \ref{bdfbv11267}), 
we know that 
\begin{align}\label{ahjd 710}
\sqrt{\Lambda e \Lambda} \otimes_{\Lambda} \Lambda u \cong 
\mathrm{Span}_{k}\Big\{ e' \otimes f \in 
\Lambda u \big\vert \exists m \in 
\mathbb{N}\mbox{ such that }
e' \otimes f^m \in \Lambda e \Lambda  
\Big\}. 
\end{align}
Thus to prove $f \in I_u$, it is enough to prove that 
\begin{align}\label{hxowc453}
(1 \otimes f^l)u \in \Lambda e \Lambda u,
\end{align}
for some $l \in \mathbb{N}$.


\vspace{2mm}
\noindent{\bf Step 2. We reduce the proof of
(\ref{hxowc453}) to the proof of the following (\ref{ncivfbg}).}

By the definition of $\iota_{\#}^i$, we know that 
$I^i$ is generated by elements in $k[V'^{i}]^{+}$, where $V'^{i}$ is the
kernel of the canonical projection $q^i: V \rightarrow V^i$ 
(see the definition in the proof of Lemma \ref{Inject}) and 
$k[V'^{i}]^{+}$ is given by the natural splitting 
$k[V'^{i}] \cong k[V'^{i}]^{+} \oplus k$. 
Without loss of generality, by the assumption of $f$, we let 
$$
f = \bigodot\limits_{(\chi_0, G^i) \not\sim (\chi_u, G^i)} f_i,
$$
where $\bigodot$ means 
the product of $R$, and $f_i \in I^i$ is a variable. 

To this end, consider the subgroup
$\langle\chi_{f_i}\rangle$, which 
is generated by $\chi_{f_i}$, in 
the dual group
$(G^i)^{\vee}$ of $G^i$. The following equality is straightforward: 
\begin{align}\label{BUXGECG3}
\#(\langle\chi_{f_i}\rangle) \#(G_{f_i}) = \#(G^i),
\end{align}
where $\#(-)$ is
the number of elements in group and
$G_{f_i} := \{g \in G^i \mid g(f_i) = f_i \}$. 
It follows that 
\begin{align}\label{BUXGECG322}
\#(\langle\chi_{f_i}\rangle)  = \#(G^i / G_{f_i}),
\end{align}
and hence that $\langle\chi_{f_i}\rangle 
=  (G^i / G_{f_i})^{\vee}$ in $(G^i)^{\vee}$.  

Now, from the definition of the
pair $(\chi_u, G^i) \in \tilde{G}_0$, we know
that there is no other $K^l$ such that
$K^l \subseteq G^i$ and $\chi_u \mid_{K^l}$ 
is nontrivial.
Meanwhile, we have $G_{f_i} \subsetneq G^i$ and $G_{f_i} \in \{K^l\}_l$ by the fact that $f_i$ is a variable.
Then we have that $\chi_u \mid_{G_{f_i}}$ is trivial. 
Hence, $\chi_u \in (G^i/G_{f_i})^{\vee}$.
We get that $\chi_{0} \chi_{f_i}^{m} = \chi_{f_i}^{m} 
= \chi_u$ in $(G^i/G_{f_i})^{\vee}$ for some $m \in \mathbb{N}$ by $\langle\chi_{f_i}\rangle 
=  (G^i / G_{f_i})^{\vee}$. 
Then we obtain that
$\chi_{0} \chi_{f_i}^{m} = \chi_u$ in $(G^i)^\vee$ by $(G^i/G_{f_i})^{\vee} \subseteq (G^i)^\vee$, and hence
that 
\begin{align}\label{nxje987}
(\chi_u \chi_{f_i}^{m'}, G^i) \sim (\chi_0, G^i)
\end{align}
in $\tilde{G}$ 
for some $m' \in \mathbb{N}$ such that $\chi_{f_i}^{-m} = \chi_{f_i}^{m'}$.
Let $e^i$ be the indecomposable idempotent such that
$ (1 \otimes f_i^{m}) u = e^i (1 \otimes f_i^{m})u \in \Lambda u$. 
By 
\begin{align}\label{jvxw56} 
(1 \otimes f_i^{m})u &= (1 \otimes f_i^{m}) 
(\frac{1}{|G|} \sum\limits_{g \in G} \chi_u(g) g \otimes 1)	 \nonumber \\
&= \frac{1}{|G|} \sum\limits_{g \in G} 
( \chi_u(g) g  \otimes g^{-1}(f_i^{m})) \nonumber \\
& =\displaystyle
\frac{1}{|G|} \sum\limits_{g \in G} 
( \chi_{f_i}^{-m}(g) \chi_u(g) g  \otimes f_i^{m}) \nonumber \\ 
& =(\frac{1}{|G|} \sum\limits_{g \in G} 
\chi^{m'}_{f_i}\chi_u(g)g  \otimes 1 )  (1 \otimes f_i^{m}),
\end{align}
we get that $e^i = \displaystyle
\frac{1}{|G|} 
\sum\limits_{g \in G} 
\chi^{m'}_{f_i}\chi_u(g)g  \otimes 1$. 
Next, by 
(\ref{nxje987}), 
we have that 
$$
\varsigma^{i}_{\chi_0}(e^i)  = \varsigma^{i}_{\chi_0} ( \frac{1}{|G|} 
\sum\limits_{g \in G} \chi^{m'}_{f_i}\chi_u(g)g  \otimes 1 ) \neq 0,
$$ 
where $\varsigma^i_{\chi_0}$ corresponds 
to the pair $(\chi_0, G^i)$. 
By the fact that $\varsigma^{i}_{\chi_0}
\circ \kappa^{i}_{\chi_0} = \mathrm{Id}$, and the definitions of 
$\varsigma^{i}_{\chi_0}$ and $\kappa^{i}_{\chi_0}$(see the proof of Lemma \ref{Compos}), 
it follows that $e^i$ is in the image of $\kappa^i_{\chi_0}$, where 
$\kappa^i_{\chi_0}$ corresponds to the pair $(\chi_0, G^i)$.  
Thus, $e^i$ is in the image of $\bar\kappa^i_{\chi_0}$.

In what follows, we replace $G$, $V$ 
and $u$ by $G/G^i$, $V^i$ and $e^i$ 
respectively. Let $f'$ be a factor of $f$ 
satisfying 
that it is the product of 
all variables like $f_r \in \{ f_i \}
$ 
such that $G^i \subsetneq G^r $, 
where $G^r$ corresponds to $f_r$. 
We can write $f'$ as
$$
\bigodot\limits_{f_r \in k[V^i]} f_r. 
$$
Then we have 
$$ f' \in \bigcap\limits_{(\chi_0, G^r/G^i) \not\sim (\chi_{e^i}, G^r/G^i)}  I^r,
$$ 
where $\chi_{e^i}$ is the character corresponding to the
idempotent $e^i$. Now, we also replace $f$ by $f'$. 
Note that since the action of $G^i$ on each $f_r$ is trivial, 
the action of $G^i$ on $f'$ is also trivial.
It implies that for any idempotent $u', u''$ in $\mathrm{Im}(\bar{\kappa}^i_{\chi_0})$, 
$u'(1 \otimes f')u''$ is also in $\mathrm{Im}(\bar{\kappa}^i_{\chi_0})$. 


Since $(e^i \otimes f_i^{m}) \in \Lambda u $, to prove that 
$(1 \otimes f^l)u \in \Lambda e \Lambda u$
(see \ref{hxowc453}) for some 
$l \in \mathbb{N}$, 
it suffices to show that
$(1 \otimes f'^{l'})e^i \in \Lambda e \Lambda e^i$, for some
$l' \in \mathbb{N}$. 
Since $e, e^i, (1 \otimes f'^{l'})e^i \in 
\mathrm{Im}(\bar{\kappa}^i_{\chi_0})$, it is left to show that 
\begin{align}\label{ncivfbg}
	(1 \otimes f'^{l'})e^i \in \Lambda_i e \Lambda_i e^{i},
\end{align}
 where 
we identity $\Lambda_i$ with the image 
%$\mathrm{Im}(\bar{\kappa}^i_{\chi_0})$ 
of injection $\bar{\kappa}^i_{\chi_0}$. 
In other words,
the proof of the lemma for $\Lambda$ to 
is reduced to the one for $\Lambda_i$. 

\vspace{2mm}
\noindent{\bf Step 3. We prove
(\ref{ncivfbg}) by reducing it to the case that $\tilde{G}_0$ is empty.}

We repeatedly apply the above procedure. 
Note that the number of equivalence classes
of $\{\chi, G^r\}_r$ in $\tilde{G}_0$ 
such that $G^i \subsetneq G^r$
is less than the number of equivalence
classes of $\{\chi, G^i\}_i$ in 
$\tilde{G}_0$, and thus
by reduction, the proof of this lemma for $\Lambda_i$ 
is reduced to the case 
where $\Lambda$ satisfies that
the set $\tilde{G}_0$ of equivalence classes of $\{\chi, G^i\}_i$ is empty. 

In this case, we have $\chi(K^l) = 1$ for any character $\chi$ and $K^l$. Since $G \in \{K^l\}_l$, $G$ is trivial. Hence 
$(1 \otimes f'^{l'})e^i \in \Lambda_i e \Lambda_i e^{i}$(see \ref{ncivfbg}) holds immediately. 
The proof is now complete. 
%
%which implies that 
%$\chi(H^j) = 1$ for any $\chi$ and $H^j$. 
%
%We claim that $H^j$ is a trivial
%subgroup. Otherwise, if not,
%then 
%there exists a character $\chi$,  different
%from $\chi_0$, 
%such that $\chi(H^j) \neq 1$. 
%Since the action of $G$ on $V$ is faithful, 
%there is a nontrivial element $h \in R$ that 
%satisfies $h \notin R^{H^j}$,
%and then $\chi_{h}(H^j) \neq 1$, which is a
%contradiction.
%
%We thus have $\{ H^j \}_j = \{1_G\}$. 
%It follows
%that for
%any nontrivial subgroup $H \subseteq G$,
%its invariant subspace $V^H \subseteq V$ is trivial. %Hence for any
%variable $x_j \in R$, we have $G_{x_j} \cong 1_G$, where %$G_{x_j} := \{g \in G\mid g(x_j) = x_j \}$. 
%In the meantime, we know that
%\begin{align*}%\label{1ne8y3}
%\#(\langle\chi_{x_j}\rangle) \#(G_{x_j}) = \#(G),
%\end{align*}
%which implies that $\langle\chi_{x_j}\rangle \cong %G^{\vee}$. 
%This means that $\chi_{x_j}$ generates $G^\vee$. 
%Thus we get that 
%$$
%\chi_0 \chi_{x_j}^{s} = \chi_{x_j}^{s} = \chi_u,
%$$ 
%for some $s \in \mathbb{N}$. 
%By the formula similar to
%\eqref{jvxw56}, we have that
%$$
%e (1 \otimes \chi_{x_j}^s) u =  (1 \otimes \chi_{x_j}^s) u,
%$$
%and then 
%$(1 \otimes x_j^s) u \in \Lambda e \Lambda u$. 
%Since $x_j$ is arbitrary, 
%we may 
%set $x_j$ be a factor of $f$,
%and then we get
%$(1 \otimes f^s) u \in \Lambda e \Lambda u$ as required by %\eqref{hxowc453}. 
%We thus have completed the proof. 
\end{proof}

\end{document}
